\begin{prob}
\textbf{Lab1 Q1: Return and Print.} Predict the full prompt display for \py{welcome()}, then for \py{print(welcome(),cal())}.
\codeblock{lab_print}
\emph{Hint:} Evaluate the two operands left to right; record each printed line separately from the strings returned.
\end{prob}
\begin{prob}
\textbf{Lab1 Q2: WWPD: What If?} Predict the full prompt display for \py{ab(10,20)}, \py{bake(0,5)}, \py{bake(1,'yum')}, and \py{bake(2,3)}.
\codeblock{lab_if}
\emph{Hint:} The two \py{if}s in \py{bake} are separate; the \py{else} belongs to the second. Printing does not exit a call.
\end{prob}
\begin{prob}
\textbf{Lab1 Q3: Pick a Digit.} For positive integer $ n $ and nonnegative integer $ k $, return the digit $ k $ places left of the rightmost digit, or $0$ if absent. Use a single return statement. Tests: $(3579,2)\mapsto5$, $(3579,0)\mapsto9$, $(3579,10)\mapsto0$.\par
\emph{Hint:} Remove $ k $ rightmost digits using division by $10^k $, then extract one digit with remainder.
\end{prob}
\begin{prob}
\textbf{Lab1 Q4: Middle Number.} Return the middle of three distinct integers using one return expression. Tests $(3,5,4)\mapsto4$, $(30,5,4)\mapsto5$, $(3,5,40)\mapsto5$, $(30,5,40)\mapsto30$.\par
\emph{Hint:} Sum all three, then remove the minimum and maximum; arithmetic is exact for these integers.
\end{prob}
\begin{prob}
\textbf{Lab1 Q5: Falling Factorial.} Return $ n(n-1)\cdots(n-k+1)$, where $ k $ is a nonnegative integer counting the consecutive descending factors; for $ k=0$, return $1$. Tests $(6,3)\mapsto120$, $(4,3)\mapsto24$, $(4,1)\mapsto4$, $(4,0)\mapsto1$.\par
\emph{Hint:} Keep the product of the factors already used and the number of factors remaining.
\end{prob}
\begin{prob}
\textbf{Lab1 Q6: Divisible By k.} For positive integers $ n,k $, print all positive multiples of $ k $ at most $ n $, in increasing order, and return their count. Test $(10,2)$: print $2,4,6,8,10$, return $5$; $(3,1)$ returns $3$; $(6,7)$ prints nothing and returns $0$.\par
\emph{Hint:} Start at $ k $ and advance by $ k $; a separate counter records printed lines.
\end{prob}
\begin{prob}
\textbf{Lab1 Q7: Double Eights.} For a nonnegative integer, determine whether the decimal digits contain adjacent $8$ s. Source true cases: $88,2882,880088$; false: $8,12345,80808080,1077$.\par
\emph{Hint:} Check the last two digits, then remove only the last digit.
\end{prob}
\begin{prob}
\textbf{Lab1 Q8: Syllabus Quiz.} The \href{https://docs.google.com/forms/d/e/1FAIpQLSdF_ze5dj1Iy5AyqCFkgk37dbThV8GnYqC7JI4e5tOH_X4RWQ/viewform}{official Syllabus Quiz} is a course-policy form that requests sign-in. Its question text is unavailable in the public materials.\par
\emph{Hint:} Use the actual linked form and official syllabus when reviewing that administrative requirement.
\end{prob}
\begin{prob}
\textbf{Lab1 Q9 (optional): Sum Digits.} Sum the decimal digits of a nonnegative integer. Tests $10\mapsto1$, $4224\mapsto12$, $1234567890\mapsto45$, $123\mapsto6$.\par
\emph{Hint:} Add the last digit and remove it; before each test, the accumulator is the sum already removed.
\end{prob}
\begin{prob}
\textbf{Discussion1 Q1: Race.} The tortoise walks $ x $ feet/minute. The hare alternates five minutes at $ y $ feet/minute and five resting, with integers $0<x<y\le2x $. Find an input for which the supplied function is wrong or runs forever, and explain why. The source's intended equality examples are \py{race(5,7)}$=7$ and \py{race(2,4)}$=10$.
\codeblock{race_source}
\emph{Hint:} A nonzero signed distance difference is true on either side of zero. A first crossing may skip equality.
\end{prob}
\begin{prob}
\textbf{Discussion1 Q2: Fizzbuzz.} For each integer $1$ through positive $ n $, print \py{fizzbuzz} for multiples of both $3$ and $5$, \py{fizz} for other multiples of $3$, \py{buzz} for other multiples of $5$, otherwise the number. The source tests $ n=16$ and confirms the returned value is \py{None}.\par
\emph{Hint:} Put the joint condition first in an \py{if}/\py{elif}/\py{else} chain.
\end{prob}
\begin{prob}
\textbf{Discussion1 Q3: Is Prime?} Return whether positive integer $ n $ is prime. Tests: $10\mapsto\py{False}$, $7\mapsto\py{True}$, $1\mapsto\py{False}$. Describe the process in one sentence.\par
\emph{Hint:} Handle $ n<2$, then search for a divisor from $2$ through $ n-1$; one divisor suffices to disprove primality.
\end{prob}
\begin{prob}
\textbf{Discussion1 Q4: Unique Digits.} Implement \py{has\_digit(n,k)} for digit $0\le k<10$, and call it from \py{unique\_digits(n)} to count the distinct digits of positive $ n $. Tests: $8675309\mapsto7$, $13173131\mapsto3$, $101\mapsto2$; \py{has\_digit(10,1)} true, \py{has\_digit(12,7)} false.\par
\emph{Hint:} Check each of the ten possible digits once; each helper searches the original number independently.
\end{prob}
\begin{prob}
\textbf{Discussion1 Q5: Bottles.} For the code below, explain what determines the number of call frames, what use is made of \py{pass\_it}'s return value, and whether its local assignment changes the global frame.
\codeblock{bottles}
\emph{Hint:} Count executed user-function calls, then distinguish global \py{bottles} from the callee's local name.
\end{prob}
\begin{prob}
\textbf{Discussion1 Q6: Double Trouble.} Draw the environment diagram after this code. Which names refer to the same function? What do \py{hat(4)} and \py{double(4)} return?
\codeblock{double_trouble}
\emph{Hint:} Assignment copies the current function value, not a reference to the other variable name. The source asks for the diagram; the two calls are additional checks.
\end{prob}
\begin{prob}
\textbf{Discussion1 Q7 (optional; Spring2023 MT1 variant): Repeating.} Return whether positive integer $ n $ consists of one or more repeats of some $ t $-digit positive integer. For $ n=6161$, results for $ t=1,2,3,4,5$ are false, true, false, true, false. No \py{for}, \py{range}, lists, or string conversion.\par
\emph{Hint:} Save the last $ t $ digits and compare equal-width blocks while removing them. Also require that saved block to have $ t $ digits: the worksheet's original initial-length guard alone misses this case.
\end{prob}
\begin{prob}
\textbf{HW2 Q1: Product.} For positive integer $ n $ and one-argument \py{term}, return \py{term(1)*...*term(n)}. Supplied helpers: identity, square, increment, triple. Source tests: product of first $3$ identities $=6$, first $5$ identities $=120$, first $3$ squares $=36$, first $5$ squares $=14400$, first $3$ increments $=24$, first $3$ triples $=162$.\par
\emph{Hint:} Use the summation loop's maintained property with multiplication and initial result $1$.
\end{prob}
\begin{prob}
\textbf{HW2 Q2: Accumulate.} Combine \py{start} and the first $ n\ge0$ terms using a supplied two-argument \py{fuse}. Assume it is commutative ($ f(a,b)=f(b,a)$) and associative ($ f(f(a,b),c)=f(a,f(b,c))$). Then implement \py{summation\_using\_accumulate} and \py{product\_using\_accumulate}, each as a single return line calling \py{accumulate}.\par
\emph{Hint:} The start value is returned unchanged for zero terms; choose identity starts $0$ and $1$ for the wrappers.
\end{prob}
\begin{prob}
\textbf{HW2 Q3: Make Repeater.} Return a function applying the supplied one-argument \py{f} exactly positive integer $ n $ times. Tests: three increments at $5\mapsto8$; five triples at $1\mapsto243$; two squares at $5\mapsto625$; three squares at $5\mapsto390625$.\par
\emph{Hint:} Capture \py{f,n} in the defining frame; create the current value and counter separately for each returned-function call.
\end{prob}
\begin{prob}
\textbf{HW2 Q4: Composite Identity Function.} Return a one-argument function testing \py{f(g(x)) == g(f(x))}; each composition's intermediate value is a valid input for the other function. Square and increment agree at $0$, but not at $4$.\par
\emph{Hint:} Return the comparison function; do not perform a comparison before its input $ x $ is supplied.
\end{prob}
\begin{prob}
\textbf{HW2 Q5 (just for fun, outside required course scope): Church Numerals.} Given \py{zero} and \py{successor}, implement \py{one} and \py{two} directly without calling \py{successor}; then implement conversion \py{church\_to\_int}, addition \py{add\_church}, multiplication \py{mul\_church}, and power \py{pow\_church}. Source checks convert zero/one/two/three to $0/1/2/3$, $2+3$ to $5$, $2\cdot3$ to $6$, $3\cdot4$ to $12$, $2^3$ to $8$, $3^2$ to $9$.\par
\emph{Hint:} Reason about how many applications of an arbitrary \py{f} each returned function performs.
\end{prob}

\subsection{Hog project: Phase1 due September10, complete project September17}
\begin{prob}
\textbf{Hog P0: Dice.} Predict five successive calls to \py{make\_test\_dice(4,1,2)}'s returned die. Explain why it is not pure.\par
\emph{Hint:} The supplied die retains its position in a repeating cycle.
\end{prob}
\begin{prob}
\textbf{Hog P1: Roll Dice.} Implement \py{roll\_dice(num\_rolls,dice=six\_sided)} for positive roll count, with Sow Sad. Call the die exactly \py{num\_rolls} times, including after a $1$. No lists, strings, square brackets, or \py{for} statements; preserve the supplied assertions.\par
\emph{Hint:} Keep both a sum and a Boolean recording whether any $1$ has occurred; choose the return value only after all calls.
\end{prob}
\begin{prob}
\textbf{Hog P2: Boar Brawl.} Return the zero-roll award from own and opponent scores; handle arbitrary nonnegative scores, including above $99$. No lists, strings, square brackets, or \py{for} statements.\par
\emph{Hint:} Extract own ones and opponent tens with remainder; award at least $1$.
\end{prob}
\begin{prob}
\textbf{Hog P3: Take Turn.} Implement \py{take\_turn} for roll counts $0$--$10$, calling \py{boar\_brawl} or \py{roll\_dice} instead of repeating their implementations; preserve supplied assertions.\par
\emph{Hint:} Dispatch on zero; pass the supplied die through only to the positive-roll case.
\end{prob}
\begin{prob}
\textbf{Hog P4: Sus Fuss.} Implement \py{num\_factors}, \py{sus\_points}, and \py{sus\_update}. Use the supplied \py{is\_prime} and the earlier functions. \py{sus\_points} returns an updated total, unchanged unless it has $3$ or $4$ factors.\par
\emph{Hint:} Count inclusive divisors; find the next prime only for a sus total; apply the rule after adding the turn's points.
\end{prob}
\begin{prob}
\textbf{Hog P5: Play.} Simulate alternating turns with supplied strategies and update, optional starting scores/die/goal. Call only the current player's strategy, once per turn; return final totals in Player0, Player1 order.\par
\emph{Hint:} Store whose turn is next and reverse the score arguments for Player1. An update returns a total, not a gain.
\end{prob}
\begin{prob}
\textbf{Hog P6: Always Roll.} Return a strategy that accepts two scores and always returns the captured $0\le n\le10$.\par
\emph{Hint:} The outer function chooses $ n $; the returned function accepts scores later.
\end{prob}
\begin{prob}
\textbf{Hog P7: Is Always Roll.} Check whether a deterministic strategy returns the same roll count for every pair of scores below the specified positive goal.\par
\emph{Hint:} Two nested loops cover every pair, including unequal scores; stop at the first mismatch.
\end{prob}
\begin{prob}
\textbf{Hog P8: Make Averaged.} Return a function accepting the same arguments as the supplied function, calling it \py{iterations} times on those arguments, and returning the mean. Default iterations $=1000$; use \py{*args}.\par
\emph{Hint:} Capture the function and positive iteration count; create fresh sum and counter on each returned-function call.
\end{prob}
\begin{prob}
\textbf{Hog P9: Max Scoring Num Rolls.} Use \py{make\_averaged} and \py{roll\_dice} to compare roll counts $1$--$10$ in ascending order, returning the smaller count on equal sampled means.\par
\emph{Hint:} Replace the current best only on a strict improvement.
\end{prob}
\begin{prob}
\textbf{Hog P10: Boar Strategy.} Return zero rolls if Boar Brawl gives at least the threshold, otherwise \py{num\_rolls}. Ignore Sus Fuss. Defaults: threshold $=11$, fallback $=6$.\par
\emph{Hint:} Compare the direct zero-roll award, not a new total.
\end{prob}
\begin{prob}
\textbf{Hog P11: Sus Strategy.} Same interface/defaults as P10, but compare the gain after both zero-roll and Sus Fuss rules with the threshold.\par
\emph{Hint:} Subtract the starting score from \py{sus\_update(0,score,opponent)}.
\end{prob}
\begin{prob}
\textbf{Hog P12 (optional): Final Strategy.} Supply a valid strategy using the rules to improve play against the baseline. There is no unique required algorithm.\par
\emph{Hint:} A zero-roll turn that certainly reaches the goal already wins; otherwise choose a justified baseline rule and test it.
\end{prob}

\subsection{Actual past-exam practice}
\begin{prob}
\textbf{Fall2022 Midterm1 Q1(c), reproduced in Lecture5 slide4.} Predict every printed line of \py{print(print(print(s,s) or print("Who's There?")), "Who?")} with \py{s='Knock'}.\par
\emph{Hint:} Each inner print returns false-valued \py{None}; \py{or} therefore evaluates its right operand.
\end{prob}
\begin{prob}
\textbf{Fall2022 Midterm1 Q3(a), reproduced in Lecture5 slides6--7.} For integer $ n>1$, sum its proper factors and return \py{'deficient'}, \py{'perfect'}, or \py{'abundant'} when that sum is respectively below, equal to, or above $ n $. Tests: $6$ perfect, $24$ abundant, $23$ deficient.\par
\emph{Hint:} Sum divisors strictly below $ n $, then compare the resulting total with $ n $.
\end{prob}
\begin{prob}
\textbf{Fall2026 Midterm1 Q3: Long Odds, p.4.} A group is a maximal consecutive sequence of all-odd or all-even decimal digits. Count the groups of positive $ n $ with length at least positive $ k $. For $15732429999$, the groups are $1573,242,9999$; at thresholds $3,4,5,1$ the counts are $3,2,0,3$.

The supplied loop increases \py{group\_size} before each digit removal, increments \py{count} at condition(a), resets \py{group\_size=0} at condition(b), then sets \py{n=n//10}. Choose(a) from \py{group\_size < k}, \py{<= k}, \py{> k}, \py{>= k}, \py{== k}, \py{!= k}, \py{\% k == 0}, \py{\% k != 0}; fill(b). In the two answers, do not write \py{str}, \py{len}, \py{if}, \py{else}, square brackets, or digits $3$--$9$.\par
\emph{Hint:} Count a group exactly when it first reaches the threshold; a parity change separates groups. The given loop already contains its two \py{if} statements.
\codeblock{long_groups_stub}
\end{prob}
\begin{prob}
\textbf{Fall2026 Midterm1 Q4.1--Q4.2: Space Between, p.5.} For $ n>1$, implement \py{gap(n,term)}, the largest absolute difference between consecutive values \py{term(1),...,term(n)}. In the supplied loop beginning \py{k,largest=1,0}, choose the condition from \py{k < n-1}, \py{k < n}, \py{k < n+1}, \py{k < n+2}, and fill the updated \py{largest}. Then implement \py{odd\_gap(n,term)} for $ n>2$ by filling \py{return gap((n+1)//2, \_\_\_\_\_)} so it compares only indices $1,3,5,\ldots\le n $. Source checks: \py{gap(5,square)}$=9$, \py{gap(6,lambda x:x\%4)}$=3$, \py{gap(4,lambda x:-x)}$=1$, \py{odd\_gap(5,square)}$=16$, \py{odd\_gap(8,lambda x:x\%5)}$=3$. No \py{if}, \py{else}, square brackets, \py{for}, or \py{range} in the answers.\par
\emph{Hint:} Visit adjacent index pairs $(k,k+1)$; for odd indices, translate index $ j $ to $2j-1$ inside a function.
\codeblock{gap_stub}
\end{prob}
